\documentclass[12pt]{article}
\usepackage[utf8]{inputenc}
\usepackage[T1]{fontenc}
\usepackage[margin=1in]{geometry}
\usepackage{setspace}
\usepackage{xcolor}
\usepackage[colorlinks=true, urlcolor=blue!60!black, linkcolor=blue!60!black]{hyperref}
\setlength{\parindent}{2em}
\setlength{\parskip}{1.25em}
\renewcommand{\baselinestretch}{1.0}
\begin{document}
\begin{center}
{\LARGE \textit{Prigg v. Pennsylvania}, 41 U.S. 539 (1842)\par}
\vspace{0.8em}
{\large Case Brief --- Four-Page Edition\par}
\end{center}
\vspace{1em}

\section{Memory Jogger}

The Court struck down Pennsylvania's anti-kidnapping statute as violating the Fugitive Slave Clause and the Fugitive Slave Act of 1793, holding that Congress's power over fugitive-slave rendition is exclusive, while indicating that States cannot be compelled to assist federal enforcement (\href{https://supreme.justia.com/cases/federal/us/41/539/}{Justia}).

\section{Detailed Case Facts}

Margaret Morgan was held as a slave in Maryland but had lived in virtual freedom without formal emancipation. In 1832 she moved to York County, Pennsylvania, where she gave birth to at least one child more than a year after leaving Maryland. In 1837 Margaret Ashmore appointed Edward Prigg to return Morgan. Prigg obtained a warrant from a local justice of the peace, but when the justice refused to take further cognizance of the matter, Prigg and three companions forcibly carried Morgan and her children to Maryland without any judicial certification.

The Fugitive Slave Clause provided that a person ``held to Service or Labour in one State'' who escaped into another ``shall be delivered up on Claim of the Party to whom such Service or Labour may be due,'' and the 1793 Act let claimants seize fugitives and seek removal certificates from federal judges or state magistrates. Pennsylvania's 1826 statute made it a felony, punishable by seven to twenty-one years at hard labor, to forcibly remove any Black person from the State with intent to enslave, and required claimants to use specified state procedures.

\section{Procedural History}

Prigg was indicted for kidnapping. Maryland and Pennsylvania arranged a friendly test case: an agreed special verdict, a conviction, and a \textit{pro forma} affirmance by the Pennsylvania Supreme Court. Prigg brought a writ of error under Section 25 of the Judiciary Act of 1789; the case was argued in February 1842 and decided on March 1, 1842.

\section{Judicial Votes}

Reversed, \textbf{8--1}, with seven separate opinions.
\begin{itemize}
\item \textbf{Opinion of the Court:} Justice Story, joined by Justices Catron, McKinley, and Wayne (who also concurred separately).
\item \textbf{Concurring in the judgment, rejecting exclusivity:} Chief Justice Taney, Justice Thompson, and Justice Daniel.
\item \textbf{Concurring in the judgment:} Justice Baldwin (the Clause is wholly self-executing).
\item \textbf{Dissenting:} Justice McLean.
\end{itemize}
Six Justices supported exclusive federal power; Taney, Thompson, and Daniel would have let States legislate in aid of slaveholders.

\section{Holding}

The Fugitive Slave Clause confers a positive, self-executing right to peaceably recapture a fugitive slave in any State, so any state law that interrupts, delays, or punishes recaption is void. Congress has implied power to enact the 1793 Act, and that power is exclusive, preempting all state legislation on the subject. States cannot be compelled to enforce the federal Act, and state legislatures may forbid their magistrates to act under it. Taney and Daniel objected that exclusivity and the non-compulsion language went beyond the narrow conflict before the Court, but both proved enormously consequential.

\section{Analysis of Opinions}

\textbf{Story (for the Court).} The Clause was a ``fundamental article without the adoption of which the Union could not have been formed.'' Under \textit{Somerset's Case}, free States would otherwise owe slave status no recognition, but the Constitution displaced comity with an unqualified right, and any state law that ``interrupts, limits, delays, or postpones'' possession operates as a forbidden discharge. Because ``where the end is required, the means are given,'' Congress may legislate, and under \textit{Houston v. Moore} and \textit{Sturgis v. Crowninshield} its power is exclusive, barring even ``auxiliary'' state laws. States ``cannot, therefore, be compelled to enforce'' the Act, but retain a police power over runaways. The opinion ignored the finding that one child was born free in Pennsylvania.

\textbf{Wayne (concurring).} The Clause was an indispensable compromise that neither side could modify through state remedies.

\textbf{Taney (concurring in part, dissenting in part).} The Clause forbids only state laws that \textit{impair} the right, not laws that \textit{support} it; with few federal judges, exclusivity would make the 1793 Act ``ineffectual and delusive.''

\textbf{Thompson, Baldwin, and Daniel.} Thompson rested on conflict with the 1793 Act; Baldwin found no kidnapping of an admitted slave; Daniel invoked Commerce Clause precedents to defend concurrent state power.

\textbf{McLean (dissenting).} He agreed that federal power was exclusive but rejected forcible private recaption: a free State presumes everyone free, so Prigg had to bring Morgan before a federal judge rather than abduct her.

\section{Examples: Future Applications}

\begin{enumerate}
\item \textbf{Same-Side (Preempted): Extra Rail Certificates.} A State adds inspection certificates and waiting periods to a comprehensive federal railroad-safety scheme; under \textit{Prigg}'s exclusivity logic, the ``auxiliary'' rule is void.
\item \textbf{Same-Side (Preempted): Criminalizing Federal Transport.} A State prosecutes federal immigration agents who move detainees across state lines without a state removal certificate; it may not interpose its procedures.
\item \textbf{Opposite-Side (Valid): Withdrawing State Assistance.} A State forbids its officers to spend resources enforcing a federal program, without obstructing federal agents---valid under the ``unless prohibited by state legislation'' principle.
\item \textbf{Opposite-Side (Valid): Breach-of-the-Peace Law.} A neutral ban on reckless force in vehicle repossessions fits \textit{Prigg}'s limitation of self-help to recaption without breach of the peace.
\item \textbf{Fence-Sitter: Cooperative Enforcement.} A State copies the federal standards for hazardous materials and adds state penalties ``to aid'' federal enforcement: void under Story's exclusivity, perhaps valid under Taney's and Daniel's concurrent-power view.
\end{enumerate}

\section{Critique}

\textbf{Text and history.} Critics note that the Clause grants Congress no express legislative power and does not forbid state laws aiding rendition; Taney and Daniel showed that Story inferred both an unenumerated power and its exclusivity from silence. Don Fehrenbacher questioned Story's claim that the Union could not have formed without the Clause, which was added late and with little debate.

\textbf{Freedom or slavery?} Paul Finkelman and Barbara Holden-Smith reject William Wetmore Story's defense of his father's opinion as a covert ``triumph of freedom,'' showing that Story stripped Northern States of protections against kidnapping and privately recommended the federal commissioners later adopted in the Fugitive Slave Act of 1850.

\section{Key Quotations}

\begin{quote}
``Upon this ground, we have not the slightest hesitation in holding that, under and in virtue of the Constitution, the owner of a slave is clothed with entire authority, in every State in the Union, to seize and recapture his slave whenever he can do it without any breach of the peace or any illegal violence. In this sense and to this extent, this clause of the Constitution may properly be said to execute itself, and to require no aid from legislation, state or national.''
\end{quote}
\noindent --- Justice Story, for the Court, 41 U.S. at 613.

\begin{quote}
``The clause manifestly contemplates the existence of a positive, unqualified right on the part of the owner of the slave which no state law or regulation can in any way qualify, regulate, control, or restrain. The slave is not to be discharged from service or labor in consequence of any state law or regulation.''
\end{quote}
\noindent --- Justice Story, for the Court, 41 U.S. at 612.

\begin{quote}
``As to the authority so conferred upon state magistrates, while a difference of opinion has existed, and may exist still, on the point in different States, whether state magistrates are bound to act under it, none is entertained by this Court that state magistrates may, if they choose, exercise that authority unless prohibited by state legislation.''
\end{quote}
\noindent --- Justice Story, for the Court, 41 U.S. at 622.

\end{document}
